\begin{afterword}
\summaryhead

The post answers a reader's request, from April 2007, to explain why ``reductionistic materialism'' is obviously correct when every past civilization's philosophy is suspect. Leaving materialism for later, the author holds that reductionism, in the sense to be given, is obviously correct. Science does not go back to theories as refuted as Newtonian mechanics, and reductionism is less a positive belief than the absence of one.

A Navy gunner once claimed that Newtonian mechanics, not relativity, governs artillery shells; the author replied that relativity always gives the more accurate answer. Physics uses one fundamental theory for a 747 and for colliding gold nuclei, but different models, because a fundamental model of the airplane would be impossibly costly to compute. A fundamental model would contain all the facts about wings and lift implicitly. The sense that the airplane is made of interacting levels is how an efficient multi-level model feels from inside. Reductionism is the disbelief that the higher levels of our models exist in the territory.

\responsehead

The post states a familiar position, and its central distinction is right. Using many models for one airplane, because the fundamental model cannot be computed, is a fact about our maps. Anderson made this distinction in ``More Is Different'' (1972): reducing everything to fundamental laws ``does not imply the ability to start from those laws and reconstruct the universe.'' Philosophers call the post's thesis ontological reduction, as against reduction between theories or models. The post names neither, but its thesis is the one they describe, and Anderson also reports that most working scientists accept it.

The case for ``obviously correct'' is thinner than the claim. The Newtonian comparison supports confidence ``when'' a rival view has been falsified, and the post does not say which view about reduction was falsified, or how. It then calls reductionism an absence of belief, but later states the same thesis positively: ``there is only the most basic level.'' The two framings are one claim.

The opposition the post answers is weak. It consists of a gunner who makes an error in physics and an imagined antireductionist who speaks one line. The post does not take up the philosophers who argue, from multiple realizability, that higher-level properties are distinct from physical ones even though every particular thing is physical.

The request the post answers reads, in full, as sarcasm. Matthew C's comment ends by proposing ``well-placed rocks thrown at anyone who would challenge the current paradigm''. The post's excerpt marks its cuts but leaves out the parts that show the tone.

One slip in the physics recurs. The post says the airplane obeys ``chromodynamics'' and uses ``chromodynamic'' three more times for the fundamental model. Chromodynamics is the theory of the strong force; the airplane's aerodynamics, like ordinary matter generally, is governed by electromagnetism, and the airplane contains electrons as well as quarks. The argument survives with the Standard Model in place of chromodynamics.

In short: a clear statement of ontological reductionism, close to what Anderson set out in 1972, called obviously correct and argued against a mistaken gunner and a one-line imagined opponent rather than its serious critics.
\end{afterword}
